\title{Επικουρικές Υπηρεσίες Σχετικές με τη Συχνότητα}

\section{Εισαγωγή}
\label{kef-03-theoria}

Η συχνότητα ενός συστήματος ηλεκτρικής ενέργειας παραμένει σταθερή μόνο όσο η παραγωγή ισούται
σε κάθε στιγμή με τη ζήτηση και τις απώλειες. Πιθανή ανισορροπία μεταξύ παραγωγής και κατανάλωσης που απομένει μετά τον
προγραμματισμό της λειτουργίας και τις αγορές (Κεφάλαιο 1, ενότητα «\nameref{sec:stages}»), λόγω, π.χ., σφαλμάτων πρόβλεψης, καλύπτονται από τις \emph{επικουρικές
υπηρεσίες σχετικές με τη συχνότητα}. Ο ΔΣΜΗΕ δεσμεύει τις αντίστοιχες εφεδρείες ενεργής ισχύος από τους παρόχους υπηρεσιών εξισορρόπησης και τις
ενεργοποιεί αυτόματα ή με εντολές κατανομής.

Κάθε \emph{σύγχρονη περιοχή} (synchronous area) λειτουργεί με κοινή συχνότητα. Στην Ευρώπη οι
κύριες σύγχρονες περιοχές είναι τέσσερις (Σχήμα~\ref{fig:sync-areas}): η Ηπειρωτική Ευρώπη, στην οποία ανήκει το
ΕΣΜΗΕ, η Βόρεια Ευρώπη (Nordic), με τη Νορβηγία, τη Σουηδία, τη Φινλανδία και την ανατολική
Δανία, η Μεγάλη Βρετανία και η Ιρλανδία μαζί με τη Βόρεια Ιρλανδία. Μικρότερες
σύγχρονες περιοχές σχηματίζουν τα συστήματα της Κύπρου, της Ισλανδίας και της Κορσικής με
τη Σαρδηνία. Οι σύγχρονες περιοχές συνδέονται μεταξύ τους μόνο με συνδέσεις HVDC, οπότε κάθε περιοχή ρυθμίζει τη συχνότητά της χωριστά.

\begin{figure}[!htbp]
\centering
\includegraphics[width=0.8\linewidth]{figures/sync-areas.svg}
\caption{Σύγχρονες περιοχές της Ευρώπης (2026). Η δυτική Δανία ανήκει στην Ηπειρωτική Ευρώπη
και η ανατολική στη Βόρεια Ευρώπη. Σύγχρονα με την Ηπειρωτική Ευρώπη λειτουργούν και τα
συστήματα της Ουκρανίας, της Μολδαβίας, της Τουρκίας, του Μαρόκου, της Αλγερίας και της
Τυνησίας. Με το χρώμα των μικρότερων περιοχών σημειώνονται και τα νησιωτικά συστήματα που δεν
συνδέονται με εναλλασσόμενο ρεύμα με την ηπειρωτική χώρα τους, όπως οι Βαλεαρίδες Νήσοι, η
Γκότλαντ και τα μη διασυνδεδεμένα νησιά του Αιγαίου· με γκρι τα συστήματα εκτός των
παραπάνω περιοχών. Τα όρια ακολουθούν κατά προσέγγιση τα σύνορα των κρατών. Γεωγραφικά
δεδομένα: \href{https://www.naturalearthdata.com}{Natural Earth}.}
\label{fig:sync-areas}
\end{figure}

Το κεφάλαιο εξετάζει αρχικά τις κατηγορίες εφεδρειών συχνότητας και στη συνέχεια τους
μηχανισμούς στους οποίους στηρίζονται: την αδράνεια των σύγχρονων μηχανών και την εικονική
αδράνεια των μονάδων που συνδέονται μέσω μετατροπέων, την εφεδρεία διατήρησης συχνότητας και την εφεδρεία αποκατάστασης συχνότητας, η οποία
υλοποιείται από την ΑΡΠ, σε απομονωμένο και σε διασυνδεδεμένο σύστημα. Κλείνει με τις παραμέτρους συχνότητας της
Ηπειρωτικής Ευρώπης και με τις περιοχές συχνότητας και τάσης εντός των οποίων οι μονάδες
πρέπει να παραμένουν σε λειτουργία.

\section{Κατηγορίες εφεδρειών συχνότητας}
\label{sec:reserves}

Οι κατηγορίες εφεδρειών και οι παράμετροι ποιότητας συχνότητας κάθε σύγχρονης περιοχής
ορίζονται από τον SOGL, ενώ η δέσμευση και η ενεργοποίησή τους μέσω αγορών ρυθμίζονται από
τον Κανονισμό (ΕΕ) 2017/2195 για την εξισορρόπηση (Electricity Balancing Guideline, EBGL). Στην
Ελλάδα οι εφεδρείες δεσμεύονται και ενεργοποιούνται μέσω της Αγοράς Εξισορρόπησης του ΑΔΜΗΕ
(Κεφάλαιο 1, υποενότητα «\nameref{sec:greek-market}»).

Μετά από ανισορροπία μεταξύ παραγωγής και ζήτησης, όπως η απώλεια μονάδας ηλεκτροπαραγωγής, η
συχνότητα ελέγχεται από μια ακολουθία μηχανισμών με διαδοχικά βραδύτερη απόκριση. Αμέσως μετά
τη διαταραχή η ανισορροπία καλύπτεται από την κινητική ενέργεια των στρεφόμενων μαζών των
σύγχρονων μηχανών, χωρίς ενέργεια ελέγχου. Ακολουθούν οι εφεδρείες ενεργής ισχύος που ορίζει ο
SOGL:

\begin{itemize}
\item \textbf{Εφεδρεία διατήρησης συχνότητας} (Frequency Containment Reserve, FCR).
      Αυτόματη και τοπική, ενεργοποιείται από τους ρυθμιστές στροφών των μονάδων από κοινού σε
      ολόκληρη τη σύγχρονη περιοχή. Συγκρατεί και σταθεροποιεί τη συχνότητα, με μόνιμη
      απόκλιση ως προς την ονομαστική τιμή.
\item \textbf{Εφεδρεία αποκατάστασης συχνότητας} (Frequency Restoration Reserve, FRR).
      Επαναφέρει τη συχνότητα στην ονομαστική τιμή και τις ανταλλαγές ισχύος μεταξύ περιοχών
      στα προγράμματα και αποδεσμεύει την ΕΔΣ. Διακρίνεται σε δύο μορφές:
      \begin{itemize}
      \item \textbf{αΕΑΣ} (automatic Frequency Restoration Reserve, aFRR), με αυτόματη ενεργοποίηση: κεντρική, οδηγείται από το σφάλμα
            ελέγχου αποκατάστασης συχνότητας (ΣΕΑΣ, ενότητα «\nameref{sec:frr-multi}») μέσω της ΑΡΠ του
            ΔΣΜΗΕ·
      \item \textbf{χΕΑΣ} (manual Frequency Restoration Reserve, mFRR), με χειροκίνητη ενεργοποίηση: ενεργοποιείται με εντολή
            κατανομής του ΔΣΜΗΕ και αποδεσμεύει την αΕΑΣ σε παρατεταμένες ανισορροπίες.
      \end{itemize}
\item \textbf{Εφεδρεία αντικατάστασης} (ΕΑ· Replacement Reserve, RR). Αποκαθιστά το απαιτούμενο
      επίπεδο των ενεργοποιημένων ΕΑΣ, ώστε το σύστημα να είναι έτοιμο για επόμενη διαταραχή.
\end{itemize}


Κάθε επόμενη κατηγορία εφεδρείας ενεργοποιείται βραδύτερα από την
προηγούμενη και την αποδεσμεύει, ώστε οι γρηγορότερες εφεδρείες να είναι διαθέσιμες για νέα
διαταραχή. Οι χρόνοι πλήρους ενεργοποίησης δίνονται στον Πίνακα~\ref{tab:reserves}.

\begin{table}[!htbp]
\centering
\begin{tabular}{lllc}
\hline
Εφεδρεία & Λειτουργία & Ενεργοποίηση & Χρόνος πλήρους ενεργοποίησης \\
\hline
ΕΔΣ (FCR)   & διατήρηση      & αυτόματη, τοπική   & 30 s \\
αΕΑΣ (aFRR) & αποκατάσταση   & αυτόματη, κεντρική & 5 min \\
χΕΑΣ (mFRR) & αποκατάσταση   & χειροκίνητη        & 12,5 min \\
ΕΑ (RR)     & αντικατάσταση  & χειροκίνητη        & $>$ 15 min \\
\hline
\end{tabular}
\caption{Επίπεδα εφεδρείας συχνότητας στην Ηπειρωτική Ευρώπη, κατά τον SOGL και τα
τυποποιημένα προϊόντα εξισορρόπησης.}
\label{tab:reserves}
\end{table}

Οι επόμενες ενότητες εξετάζουν διαδοχικά την αδράνεια, την εικονική αδράνεια, την ΕΔΣ και
την ΕΑΣ σε απομονωμένο και σε διασυνδεδεμένο σύστημα. Το Σχήμα~\ref{fig:freq} συνοψίζει
την ακολουθία για το συμβάν αναφοράς της Ηπειρωτικής Ευρώπης.

\section{Αδράνεια και αρχικός ρυθμός μεταβολής συχνότητας}
\label{sec:inertia}

Σε μόνιμη κατάσταση η συχνότητα είναι σταθερή και κοινή σε ολόκληρη τη σύγχρονη
περιοχή· κάθε ανισορροπία μεταξύ παραγωγής και ζήτησης προκαλεί μεταβολή της. Αμέσως μετά από
ανισορροπία ισχύος, για παράδειγμα την απώλεια μεγάλης μονάδας ηλεκτροπαραγωγής, οι ρυθμιστές
στροφών δεν έχουν ακόμη αποκριθεί και το έλλειμμα καλύπτεται από την κινητική ενέργεια των
στρεφόμενων μαζών των σύγχρονων μηχανών. Ο ρυθμός με τον οποίο μεταβάλλεται τότε η συχνότητα
προκύπτει από την εξίσωση ταλάντωσης (swing equation) των μηχανών.

\subsection{Εξίσωση ταλάντωσης}
\label{sec:swing}

Θεωρείται σύγχρονη γεννήτρια χωρίς ρυθμιστή στροφών ή άλλο έλεγχο συχνότητας, η οποία
συνδέεται με το υπόλοιπο σύστημα μέσω γραμμής χωρίς απώλειες και αρχικά βρίσκεται σε
ισορροπία. Η κίνηση του δρομέα περιγράφεται από τον δεύτερο νόμο του Νεύτωνα για την
περιστροφική κίνηση:

\begin{equation}
J\,\frac{\mathrm{d}^2\delta_m}{\mathrm{d}t^2} = T_m - T_e
\label{eq:swing-torque}
\end{equation}

όπου $J$ η ροπή αδράνειας του δρομέα μαζί με τον στρόβιλο που τον κινεί, $\delta_m$ η γωνία
του δρομέα ως προς άξονα που περιστρέφεται με τη σύγχρονη ταχύτητα, $T_m$ η μηχανική ροπή του
στροβίλου και $T_e$ η ηλεκτρομαγνητική ροπή. Πολλαπλασιάζοντας με τη γωνιακή ταχύτητα του
δρομέα $\omega_m$ προκύπτει η αντίστοιχη σχέση ισχύος,
$J\omega_m\,\mathrm{d}^2\delta_m/\mathrm{d}t^2 = P_m - P_e$, όπου $P_m$ η μηχανική ισχύς του
στροβίλου και $P_e$ η ηλεκτρική ισχύς της γεννήτριας. Η αδράνεια της μηχανής εκφράζεται
συνήθως με τη \emph{σταθερά αδράνειας}

\begin{equation}
H = \frac{J\,\omega_{sm}^2}{2\,S_n}
\label{eq:hconst}
\end{equation}

δηλαδή με τον λόγο της κινητικής ενέργειας των στρεφόμενων μαζών στην ονομαστική γωνιακή
ταχύτητα $\omega_{sm}$ προς την ονομαστική φαινόμενη ισχύ $S_n$ της μηχανής. Η $H$ εκφράζεται
σε δευτερόλεπτα (MW·s/MVA) και ισούται με τον χρόνο για τον οποίο η αποθηκευμένη κινητική
ενέργεια θα μπορούσε να καλύψει την ονομαστική ισχύ της μηχανής. Με αντικατάσταση της $J$ από
την \eqref{eq:hconst}, διαίρεση με $S_n$, την προσέγγιση $\omega_m \approx \omega_{sm}$ και χρήση της ηλεκτρικής γωνίας προκύπτει η εξίσωση ταλάντωσης σε p.u.:

\begin{equation}
\frac{2H}{\omega_s}\,\frac{\mathrm{d}^2\delta}{\mathrm{d}t^2} = p_m - p_e
\label{eq:swing}
\end{equation}

όπου $\delta$ η γωνία φορτίου, $\omega_s$ η ονομαστική γωνιακή συχνότητα του δικτύου και
$p_m$, $p_e$ η μηχανική και η ηλεκτρική ισχύς ανά μονάδα, με βάση την $S_n$. Η γωνία φορτίου
μεταβάλλεται με ρυθμό $\mathrm{d}\delta/\mathrm{d}t = \omega - \omega_s$, όπου $\omega$ η
ηλεκτρική γωνιακή ταχύτητα του δρομέα, οπότε η δεύτερη παράγωγός της ισούται με
$\mathrm{d}\omega/\mathrm{d}t$. Με $\omega/\omega_s = f/f_0$, η \eqref{eq:swing} δίνει τον
RoCoF σε ανά μονάδα τιμές ανά δευτερόλεπτο (p.u./s):

\begin{equation}
\mathrm{RoCoF} = \frac{\mathrm{d}(f/f_0)}{\mathrm{d}t} = \frac{p_m - p_e}{2H}
\label{eq:rocof-pu}
\end{equation}

όπου $f_0$ η ονομαστική συχνότητα.

Σύμφωνα με την \eqref{eq:rocof-pu}, κάθε ανισορροπία μεταξύ παραγωγής και κατανάλωσης
μεταβάλλει τη συχνότητα: πλεόνασμα παραγωγής ($p_m > p_e$) την αυξάνει και έλλειμμα τη μειώνει,
με ρυθμό αντιστρόφως ανάλογο της σταθεράς αδράνειας. Η απόκριση είναι εγγενής: σε έλλειμμα ο
δρομέας επιβραδύνεται και αποδίδει κινητική ενέργεια στο δίκτυο χωρίς καμία ενέργεια ελέγχου.
Η αδράνεια δεν αποτρέπει τη μεταβολή της συχνότητας, αλλά την επιβραδύνει· επιθυμητός είναι
επομένως μικρός ρυθμός μεταβολής, ώστε η ΕΔΣ (ενότητα «\nameref{sec:fcr}») να έχει επαρκή
χρόνο να ενεργοποιηθεί πριν η συχνότητα φθάσει σε επικίνδυνες τιμές.

\subsection{Αδράνεια του συστήματος}
\label{sec:inertia-sys}

Σε ένα πραγματικό σύστημα συνεισφέρουν στη συνολική αδράνεια όλες οι στρεφόμενες μηχανές που
είναι συγχρονισμένες με το δίκτυο, είτε παράγουν είτε καταναλώνουν ενεργή ισχύ. Η σταθερά
αδράνειας του συστήματος $H_T$ ορίζεται με βάση το φορτίο του συστήματος $P_L$ και αποτελείται
από τη συνεισφορά των σύγχρονων γεννητριών και τη συνεισφορά του φορτίου:

\begin{equation}
H_T = H_{\mathrm{gen}} + H_{\mathrm{load}}
    = \frac{\sum_{i=1}^{N} H_i\,S_{G,i}}{P_L} + \frac{\sum_{j=1}^{M} H_j\,S_{L,j}}{P_L}
\label{eq:htotal}
\end{equation}

όπου $N$ το πλήθος των συγχρονισμένων γεννητριών, $H_i$ η σταθερά αδράνειας και $S_{G,i}$ η
ονομαστική ισχύς της γεννήτριας $i$, και $M$ το πλήθος των στρεφόμενων φορτίων, κυρίως των
κινητήρων που συνδέονται απευθείας στο δίκτυο, με σταθερά αδράνειας $H_j$ και ονομαστική ισχύ
$S_{L,j}$. Η σταθερά αδράνειας κάθε γεννήτριας περιλαμβάνει την κινητική ενέργεια του δρομέα
και του στροβίλου που τον κινεί. Ο ορισμός με βάση το φορτίο, και όχι την ισχύ των
συγχρονισμένων μονάδων, λαμβάνει υπόψη ότι μέρος του φορτίου καλύπτεται από μονάδες χωρίς
αδράνεια. Το γινόμενο $H_T P_L$ είναι η συνολική κινητική ενέργεια των στρεφόμενων μαζών, οπότε
από την \eqref{eq:rocof-pu}, με την ανισορροπία $\Delta P$ ανηγμένη στο φορτίο, ο αρχικός
ρυθμός μεταβολής της συχνότητας σε Hz/s είναι κατ' απόλυτη τιμή

\begin{equation}
\frac{\mathrm{d}f}{\mathrm{d}t} = \frac{f_0 \, \Delta P}{2 H_T P_L}
\label{eq:rocof}
\end{equation}

\emph{Παράδειγμα.} Για σύγχρονη περιοχή με φορτίο $P_L = 300$ GW, $H_T = 5$ s και απώλεια
$\Delta P = 3\,000$ MW, ο αρχικός ρυθμός μεταβολής είναι
$50 \times 3000 / (2 \times 5 \times 300\,000)$ = 0,05 Hz/s. Για την ίδια απώλεια σε συνθήκες
χαμηλού φορτίου, με τη μισή κινητική ενέργεια $H_T P_L$, ο ρυθμός διπλασιάζεται.

Η συνεισφορά των σύγχρονων γεννητριών εκτιμάται με σχετική ακρίβεια, επειδή τα στοιχεία των
συμβατικών σταθμών, όπως η τεχνολογία και η εγκατεστημένη ισχύς, είναι διαθέσιμα και η
λειτουργία τους καθορίζεται από τον προγραμματισμό της επόμενης ημέρας. Ο
Πίνακας~\ref{tab:hvalues} δίνει τυπικές τιμές της σταθεράς αδράνειας ανά τεχνολογία σταθμού.
Τα αιολικά και φωτοβολταϊκά πάρκα δεν συνεισφέρουν αδράνεια: τα φωτοβολταϊκά δεν διαθέτουν
στρεφόμενες μάζες, ενώ στις ανεμογεννήτριες ο μετατροπέας απομονώνει την ταχύτητα του δρομέα
από τη συχνότητα του δικτύου. Η συνεισφορά του φορτίου εκτιμάται δυσκολότερα, επειδή δεν
υπάρχει πρόσβαση σε αναλυτικά στοιχεία λειτουργίας κάθε εγκατάστασης. Η εκτίμηση από το μείγμα
παραγωγής, η οποία την αγνοεί, είναι επομένως συντηρητική. Τη συνολική αδράνεια $H_T$, μαζί με
τη συνεισφορά του φορτίου, εκτιμούν μέθοδοι που επεξεργάζονται μετρήσεις PMU (Κεφάλαιο 1,
ενότητα «\nameref{sec:wams}»).

\begin{table}[!htbp]
\centering
\begin{tabular}{lc}
\hline
Τεχνολογία σταθμού & $H$ (s) \\
\hline
Πυρηνικός                    & 5,9 \\
Πετρελαϊκός                  & 4,3 \\
Λιθανθρακικός                & 4,2 \\
Φυσικού αερίου               & 4,2 \\
Λιγνιτικός                   & 3,8 \\
Υδροηλεκτρικός με ταμιευτήρα & 3,7 \\
Αντλησιοταμίευσης            & 3,5 \\
Βιομάζας                     & 3,3 \\
Υδροηλεκτρικός ροής ποταμού  & 2,7 \\
Αιολικός, φωτοβολταϊκός      & 0   \\
\hline
\end{tabular}
\caption{Τυπικές τιμές της σταθεράς αδράνειας ανά τεχνολογία σταθμού, όπως τις χρησιμοποιεί ο
ENTSO-E για την εκτίμηση της αδράνειας της Ηπειρωτικής Ευρώπης.}
\label{tab:hvalues}
\end{table}

Σύμφωνα με την \eqref{eq:rocof}, η αντικατάσταση σύγχρονων γεννητριών από μονάδες που
συνδέονται μέσω ηλεκτρονικών μετατροπέων ισχύος, όπως τα αιολικά και φωτοβολταϊκά πάρκα,
μειώνει την αποθηκευμένη κινητική ενέργεια $H_T P_L$ και αυξάνει τον ρυθμό μεταβολής της
συχνότητας. Μειώνεται έτσι ο διαθέσιμος χρόνος πριν η συχνότητα φθάσει στο κατώφλι της
αυτόματης αποσύνδεσης ζήτησης λόγω χαμηλής συχνότητας, και απαιτούνται γρηγορότερες υπηρεσίες,
όπως η εικονική αδράνεια (ενότητα «\nameref{sec:virtual-inertia}») και η γρήγορη απόκριση
συχνότητας.

Για την Ηπειρωτική Ευρώπη, ο ENTSO-E εκτίμησε την
αδράνεια του 2019 ανά ώρα από το μείγμα παραγωγής και παρουσίασε την κατανομή της ανά μήνα και
ώρα της ημέρας, με εμφανή μείωση τις μεσημεριανές ώρες λόγω της φωτοβολταϊκής παραγωγής, η
οποία αναμένεται να ενταθεί έως το 2030 (Σχήμα~\ref{fig:inertia-ce}). Το 2019 η σταθερά
αδράνειας του συστήματος κυμαινόταν κυρίως μεταξύ 4,5 και 6 s, με τις χαμηλότερες τιμές τις
μεσημεριανές ώρες, ενώ στο σενάριο του 2030 οι μεσημεριανές τιμές της άνοιξης και του
καλοκαιριού μειώνονται περίπου στα 1,5 έως 3 s.

\begin{figure}[!htbp]
\centering
\includegraphics[width=\linewidth]{figures/entsoe-inertia-2019-2030.png}
\caption{Κατανομή της σταθεράς αδράνειας $H_{\mathrm{sys}}$ της Ηπειρωτικής Ευρώπης ανά μήνα
(οριζόντιος άξονας) και ώρα της ημέρας (κατακόρυφος άξονας), εκτιμώμενη από το μείγμα
παραγωγής, για το 2019 (αριστερά) και για το σενάριο Distributed Generation 2030 του δεκαετούς προγράμματος ανάπτυξης δικτύου (Ten-Year
Network Development Plan, TYNDP) 2018
(δεξιά). Πηγή: ENTSO-E,
\href{https://eepublicdownloads.entsoe.eu/clean-documents/SOC\%20documents/Inertia\%20and\%20RoCoF_v17_clean.pdf}{«Inertia
and Rate of Change of Frequency (RoCoF)»}, έκδοση 17, 2020, Σχήμα 10.}
\label{fig:inertia-ce}
\end{figure}

Το Σχήμα~\ref{fig:inertia} δείχνει την εξέλιξη σε ένα μικρότερο σύστημα, της Μεγάλης Βρετανίας,
για το οποίο ο διαχειριστής NESO (National Energy System Operator) δημοσιεύει ως ανοιχτά δεδομένα την αδράνεια ανά μισάωρο.
Μεταξύ 2019 και 2025 η μέση αδράνεια μειώθηκε από 241 σε 178 GVA·s, δηλαδή κατά 26\% περίπου,
το ποσοστό του χρόνου με αδράνεια κάτω από 150 GVA·s αυξήθηκε από 2\% σε 37\% και η ελάχιστη
τιμή μειώθηκε από 135 σε 108 GVA·s. Οι χαμηλότερες τιμές εμφανίζονται τις νυχτερινές ώρες και,
το 2025, και τις μεσημεριανές ώρες από την άνοιξη έως το φθινόπωρο, όταν μεγάλο μέρος της
ζήτησης καλύπτεται από αιολική και φωτοβολταϊκή παραγωγή.

\begin{figure}[!htbp]
\centering
\includegraphics[width=\linewidth]{figures/inertia-heatmap-gb.svg}
\caption{Αδράνεια του συστήματος της Μεγάλης Βρετανίας ανά μισάωρο (τοπική ώρα) για τα έτη 2019
και 2025. Κάθε κατακόρυφη λωρίδα αντιστοιχεί σε μία ημέρα και οι ανοιχτότεροι χρωματισμοί σε
χαμηλότερη αδράνεια. Δεδομένα:
\href{https://www.neso.energy/data-portal/system-inertia}{NESO, System Inertia}, Outturn Inertia.
Supported by National Energy SO Open Data.}
\label{fig:inertia}
\end{figure}

Το Σχήμα~\ref{fig:freq} απεικονίζει την απόκριση συχνότητας για το συμβάν αναφοράς και τις
τρεις διαδοχικές φάσεις της.

\begin{figure}[!htbp]
\centering
\includegraphics[width=\linewidth]{figures/frequency-response.svg}
\caption{Απόκριση συχνότητας σε απώλεια παραγωγής 3\,000 MW. Η αρχική κλίση καθορίζεται
αποκλειστικά από την αδράνεια· η ΕΔΣ ανακόπτει την πτώση
και σταθεροποιεί τη συχνότητα σε νέα τιμή μόνιμης κατάστασης· η ΕΑΣ εξαλείφει το μόνιμο σφάλμα και επαναφέρει τη συχνότητα στα 50 Hz. Η
διακεκομμένη καμπύλη είναι η απόκριση μόνο με ΕΔΣ.}
\label{fig:freq}
\end{figure}

\section{Εικονική αδράνεια}
\label{sec:virtual-inertia}

Η μείωση της αδράνειας μπορεί να αντισταθμιστεί από τις ίδιες τις μονάδες που συνδέονται μέσω
μετατροπέων. Η \emph{εικονική αδράνεια} (virtual inertia) υλοποιείται με τεχνική ελέγχου του
μετατροπέα που μιμείται τη συμπεριφορά της σύγχρονης μηχανής: σε κάθε μεταβολή της συχνότητας ο
μετατροπέας μεταβάλλει την ισχύ εξόδου του όπως ο δρομέας σύγχρονης γεννήτριας με σταθερά
αδράνειας $H$, σύμφωνα με την \eqref{eq:rocof-pu}. Ο έλεγχος ενσωματώνεται στους μετατροπείς των
μονάδων ΑΠΕ και των σταθμών αποθήκευσης. Συναφής είναι ο όρος \emph{συνθετική αδράνεια}
(synthetic inertia) του RfG, ο οποίος περιγράφει τη δυνατότητα μιας μονάδας πάρκου ισχύος
να αντικαθιστά την επίδραση της αδράνειας μιας
σύγχρονης μονάδας.

\subsection{Το μοντέλο synchronverter}
\label{sec:synchronverter}

Μία από τις πρώτες υλοποιήσεις της εικονικής αδράνειας είναι ο \emph{synchronverter}, στον
οποίο ο έλεγχος του μετατροπέα επιλύει σε πραγματικό χρόνο τις εξισώσεις της σύγχρονης
γεννήτριας. Στο Σχήμα~\ref{fig:synchronverter} ο μετατροπέας της μονάδας ΑΠΕ παράγει την
εσωτερική τάση $\mathbf{E}$ και συνδέεται στο σημείο κοινής σύνδεσης (Point of Common Coupling,
PCC), με τάση $V_{\mathrm{PCC}}$, μέσω φίλτρου με στοιχεία $R_c$, $L_c$ και $C$· το δίκτυο
παριστάνεται με την τάση $V_{\mathrm{grid}}$ πίσω από τα στοιχεία $R_g$ και $L_g$. Ο έλεγχος
αποτελείται από δύο βρόχους. Ο βρόχος ενεργής ισχύος αναπαράγει το μηχανικό μέρος της
μηχανής, δηλαδή την εξίσωση ταλάντωσης \eqref{eq:swing-torque} με πρόσθετο όρο απόσβεσης, και
ο βρόχος άεργης ισχύος τη διέγερση:

\begin{equation}
\ddot{\theta} = \frac{1}{J}\left[T_m - T_e - D_p\,(\dot{\theta} - \dot{\theta}_n)\right]
\label{eq:sv-swing}
\end{equation}

\begin{equation}
\dot{\Phi}_{sf} = k_i^v\,(Q_{\mathrm{set}} - Q)
\label{eq:sv-flux}
\end{equation}

Η ηλεκτρομαγνητική ροπή, η άεργη ισχύς και η εσωτερική τάση υπολογίζονται από τις σχέσεις της
σύγχρονης μηχανής με κυλινδρικό δρομέα:

\begin{align}
T_e &= \Phi_{sf}\,\langle \mathbf{I}, \sin\boldsymbol{\Theta} \rangle \\
Q &= -\dot{\theta}\,\Phi_{sf}\,\langle \mathbf{I}, \cos\boldsymbol{\Theta} \rangle \\
\mathbf{E} &= \dot{\theta}\,\Phi_{sf}\,\sin\boldsymbol{\Theta}
\label{eq:sv-out}
\end{align}

όπου $\theta$ η γωνία της εσωτερικής τάσης, $\dot{\theta}$ η εικονική γωνιακή συχνότητα,
$\dot{\theta}_n$ η ονομαστική γωνιακή συχνότητα, $J$ η εικονική ροπή αδράνειας, $D_p$ ο
συντελεστής απόσβεσης, $T_m = P_{\mathrm{set}}/\dot{\theta}_n$ η εικονική μηχανική ροπή που
αντιστοιχεί στην εντολή ενεργής ισχύος $P_{\mathrm{set}}$, $\Phi_{sf}$ η εικονική μαγνητική
ροή της διέγερσης, $k_i^v$ το ολοκληρωτικό κέρδος του βρόχου άεργης ισχύος και
$Q_{\mathrm{set}}$ η εντολή άεργης ισχύος. Τα $\mathbf{I}$ και $\mathbf{E}$ είναι τα διανύσματα
των τριών φασικών ρευμάτων εξόδου και των τριών εσωτερικών τάσεων,
$\boldsymbol{\Theta} = [\theta, \theta - 2\pi/3, \theta - 4\pi/3]^{T}$, το ημίτονο και το
συνημίτονο εφαρμόζονται σε κάθε στοιχείο του $\boldsymbol{\Theta}$ και
$\langle \cdot, \cdot \rangle$ είναι το εσωτερικό γινόμενο. Από την \eqref{eq:sv-out}, το μέτρο
της εσωτερικής τάσης είναι $\dot{\theta}\,\Phi_{sf}$ και η γωνία της $\theta$.

\begin{figure}[!htbp]
\centering
\includegraphics[width=\linewidth]{figures/synchronverter.svg}
\caption{Μονάδα ΑΠΕ με εικονική αδράνεια κατά το μοντέλο synchronverter. Αριστερά, το κύκλωμα
σύνδεσης: ο μετατροπέας παράγει την εσωτερική τάση $\mathbf{E}$ και συνδέεται στο σημείο κοινής
σύνδεσης (PCC) μέσω φίλτρου $R_c$, $L_c$, $C$· το δίκτυο παριστάνεται με την τάση
$V_{\mathrm{grid}}$ και τα στοιχεία $R_g$, $L_g$. Δεξιά, οι βρόχοι ελέγχου ενεργής και άεργης
ισχύος, κατά τις \eqref{eq:sv-swing} και \eqref{eq:sv-flux}· το μπλοκ των εξισώσεων εξόδου
υλοποιεί τις \eqref{eq:sv-out}.}
\label{fig:synchronverter}
\end{figure}

Παράμετροι του ελέγχου είναι οι $J$, $D_p$ και $k_i^v$, είσοδοι οι $P_{\mathrm{set}}$,
$Q_{\mathrm{set}}$, $\dot{\theta}_n$ και το μετρούμενο ρεύμα $\mathbf{I}$, εσωτερικές
μεταβλητές οι $\dot{\theta}$ και $\Phi_{sf}$ και έξοδος η τάση $\mathbf{E}$, την οποία παράγει
ο μετατροπέας. Η εικονική ροπή αδράνειας προκύπτει από την \eqref{eq:hconst} ως
$J = 2 H S_n / \dot{\theta}_n^{2}$, οπότε η σταθερά αδράνειας $H$ είναι παράμετρος σχεδιασμού
και δεν εξαρτάται από φυσική μάζα. Ο όρος απόσβεσης συγκρίνει την εικονική γωνιακή συχνότητα με
την ονομαστική και όχι με τη συχνότητα του δικτύου· σε μόνιμη κατάσταση μεταβάλλει επομένως την
ισχύ εξόδου αναλογικά προς την απόκλιση συχνότητας και λειτουργεί ως στατισμός (ενότητα
«\nameref{sec:fcr}»).

Ο synchronverter συμπεριφέρεται ως πηγή τάσης, όπως η σύγχρονη μηχανή, και είναι επομένως
μετατροπέας σχηματισμού δικτύου (Κεφάλαιο 2, υποενότητα «\nameref{sec:blackstart}»). Δεν
χρησιμοποιεί PLL για τη μέτρηση της γωνίας της τάσης του δικτύου,
αλλά παραμένει συγχρονισμένος με αυτό μέσω του βρόχου ενεργής ισχύος. Αν η συχνότητα του
δικτύου μειωθεί, η γωνία της τάσης του δικτύου υστερεί ως προς τη γωνία $\theta$, η ισχύς
εξόδου και η ροπή $T_e$ αυξάνονται και, σύμφωνα με την \eqref{eq:sv-swing}, η εικονική γωνιακή
συχνότητα $\dot{\theta}$ μειώνεται, ώστε να ακολουθήσει τη συχνότητα του δικτύου. Κατά τη
μεταβολή αυτή ο μετατροπέας αποδίδει πρόσθετη ισχύ, ανάλογη της σταθεράς αδράνειας και του
ρυθμού μεταβολής της συχνότητας, όπως ο δρομέας της σύγχρονης γεννήτριας.

\subsection{Ισοζύγιο ισχύος και διαστασιολόγηση της αποθήκευσης}
\label{sec:vi-storage}

Από την \eqref{eq:rocof-pu}, η ισχύς εξόδου μονάδας με εικονική αδράνεια είναι

\begin{equation}
p_e = p_m - 2H \cdot \mathrm{RoCoF}
\label{eq:pe-vi}
\end{equation}

όπου $p_m$ η ισχύς της πρωτογενούς πηγής ενέργειας, για παράδειγμα του ανέμου ή της ηλιακής
ακτινοβολίας, και $H$ η σταθερά εικονικής αδράνειας της μονάδας. Σε αύξηση της συχνότητας
($\mathrm{RoCoF} > 0$) η ισχύς εξόδου πρέπει να είναι μικρότερη από την ισχύ της πηγής
($p_e < p_m$) και σε μείωση της συχνότητας ($\mathrm{RoCoF} < 0$) μεγαλύτερη ($p_e > p_m$).
Η απόκλιση ισχύος καλύπτεται με δύο τρόπους:

\begin{itemize}
\item \textbf{Ρύθμιση της ισχύος της πηγής}, ώστε η $p_m$ να ακολουθεί την $p_e$. Για να
      διαθέτει πρόσθετη ισχύ σε μείωση της συχνότητας, η μονάδα πρέπει να λειτουργεί κάτω από
      το σημείο μέγιστης ισχύος (Maximum Power Point, MPP), με απώλεια ανανεώσιμης ενέργειας.
\item \textbf{Σύστημα αποθήκευσης}, το οποίο απορροφά ή αποδίδει κάθε διαφορά μεταξύ $p_e$ και
      $p_m$, με επιπλέον κόστος εγκατάστασης.
\end{itemize}

Στη δεύτερη περίπτωση η ισχύς που πρέπει να διαθέτει το σύστημα αποθήκευσης, η \emph{ισχύς
εικονικής αδράνειας}, προκύπτει από την \eqref{eq:pe-vi} ως διαφορά $p_{ir} = p_e - p_m$:

\begin{equation}
p_{ir} = -\frac{2H}{f_0}\,\frac{\mathrm{d}f}{\mathrm{d}t}
\label{eq:pir}
\end{equation}

Αν η διαταραχή συμβαίνει τη στιγμή $t_0$, με $f(t_0) = f_0$, και η συχνότητα φθάνει στην
ελάχιστη τιμή της (ναδίρ) τη στιγμή $t_1$, η \emph{ενέργεια εικονικής αδράνειας} είναι

\begin{equation}
e_{ir} = \int_{t_0}^{t_1} p_{ir}\,\mathrm{d}t = -\frac{2H}{f_0}\left(f(t_1) - f_0\right)
\label{eq:eir}
\end{equation}

σε p.u.·s, δηλαδή σε kJ ανά kW ονομαστικής ισχύος. Η ενέργεια εξαρτάται μόνο από την τελική
απόκλιση της συχνότητας και όχι από τη χρονική εξέλιξή της. Η μέγιστη απόκλιση
$\Delta f_{\max}$ για την οποία διαστασιολογείται η αποθήκευση είναι 2,5 Hz, αφού κάτω από τα
47,5 Hz η παραμονή των μονάδων σε λειτουργία δεν απαιτείται (ενότητα
«\nameref{sec:f-ranges}»). Επομένως, για $f_0 = 50$ Hz,

\begin{equation}
e_{ir}^{\max} = \frac{2H}{f_0}\,\Delta f_{\max} = \frac{H}{10}
\label{eq:eir-max}
\end{equation}

σε kJ/kW με την $H$ σε δευτερόλεπτα, δηλαδή 0,1 $H$ kJ/kW ή περίπου 0,0278 $H$ Wh/kW.

\emph{Παράδειγμα.} Για $H = 5$ s η μέγιστη ενέργεια είναι 0,5 kJ/kW, δηλαδή περίπου 0,14 Wh
ανά kW ή 0,14 kWh ανά MW ονομαστικής ισχύος. Η τιμή αυτή είναι το ένα δέκατο της κινητικής
ενέργειας, 5 kJ/kW, που αποθηκεύει σύγχρονη γεννήτρια με την ίδια σταθερά αδράνειας, επειδή
έως τα 47,5 Hz η γεννήτρια αποδίδει μόνο το μέρος της κινητικής της ενέργειας που αντιστοιχεί
στη μείωση της ταχύτητάς της. Η απαιτούμενη ενέργεια είναι επομένως μικρή σε σχέση με την
ισχύ: σύμφωνα με την \eqref{eq:pir}, για ρυθμό μείωσης της συχνότητας 1 Hz/s απαιτείται
$p_{ir} = 2 \times 5 \times 1/50$ = 0,2 p.u., δηλαδή το 20\% της ονομαστικής ισχύος της
μονάδας για λίγα δευτερόλεπτα.

\section{Εφεδρεία διατήρησης συχνότητας και στατισμός}
\label{sec:fcr}

Η ΕΔΣ υλοποιείται από τους ρυθμιστές στροφών (governors) των μονάδων, οι οποίοι μεταβάλλουν
την παραγόμενη ισχύ αναλογικά προς την απόκλιση συχνότητας. Η αναλογία καθορίζεται από τον
\textbf{στατισμό} (droop) $R$ της μονάδας, ο οποίος ορίζεται ως ο λόγος της ανά μονάδα
μεταβολής της συχνότητας προς την ανά μονάδα μεταβολή της ισχύος εξόδου σε μόνιμη κατάσταση,
με βάσεις την ονομαστική συχνότητα $f_0$ και την ονομαστική ισχύ $P_n$ της μονάδας:

\begin{equation}
R = -\frac{\Delta f / f_0}{\Delta P / P_n}
\label{eq:droopdef}
\end{equation}

ή ισοδύναμα

\begin{equation}
\frac{\Delta P}{P_n} = -\frac{1}{R}\,\frac{\Delta f}{f_0}
\label{eq:droop}
\end{equation}

Ο στατισμός είναι η κλίση της στατικής χαρακτηριστικής συχνότητας-ισχύος της μονάδας και
εκφράζεται συνήθως επί τοις εκατό. Στατισμός 5\% σημαίνει ότι μεταβολή της συχνότητας κατά 5\%
της ονομαστικής, δηλαδή κατά 2,5 Hz σε σύστημα 50 Hz, μεταβάλλει την ισχύ εξόδου κατά την
ονομαστική ισχύ της μονάδας· μικρότερος στατισμός σημαίνει ισχυρότερη απόκριση. Το αντίστροφο
$1/R$ είναι το κέρδος του αναλογικού ρυθμιστή. Με μηδενικό στατισμό (ισόχρονη ρύθμιση) η
μονάδα θα επανέφερε μόνη της τη συχνότητα στην ονομαστική τιμή, αλλά περισσότερες ισόχρονες
μονάδες δεν μπορούν να λειτουργούν παράλληλα, επειδή η κατανομή του φορτίου μεταξύ τους δεν
είναι καθορισμένη.

Σε μόνιμη κατάσταση όλες οι μονάδες της σύγχρονης περιοχής αντιλαμβάνονται την ίδια
απόκλιση συχνότητας, οπότε από την \eqref{eq:droop} η συνεισφορά της μονάδας $i$ είναι
$\Delta P_i = -P_{n,i}\,\Delta f / (R_i f_0)$: μονάδες με ίσο στατισμό μοιράζονται τη
μεταβολή φορτίου αναλογικά προς την ονομαστική τους ισχύ. Ο στατισμός καθορίζει επομένως τη
συμμετοχή κάθε μονάδας χωρίς ανάγκη επικοινωνίας μεταξύ τους. Στην Ηπειρωτική Ευρώπη η ΕΔΣ
ενεργοποιείται πλήρως για απόκλιση συχνότητας $\pm 200$ mHz, εντός 30 s, ενώ η συνολική νεκρή
ζώνη και αναισθησία της απόκρισης δεν επιτρέπεται να υπερβαίνει τα 10 mHz.

Τυπικές τιμές στατισμού είναι $R = 4$ έως 5\%. Για παράδειγμα, μονάδα 500 MW με $R = 5\%$,
για απόκλιση $\Delta f = -200$ mHz, δηλαδή −0,4\%, αυξάνει την παραγωγή της κατά
(0,004~/~0,05)~×~500~=~40~MW.

Το αρνητικό πρόσημο αντιστοιχεί σε αρνητική ανάδραση: μείωση της συχνότητας προκαλεί αύξηση
της παραγωγής. Επειδή ο έλεγχος είναι \emph{αναλογικός}, η πρόσθετη παραγωγή διατηρείται μόνο
όσο διατηρείται η απόκλιση συχνότητας. Για ανισορροπία $\Delta P$ η απόκλιση μόνιμης
κατάστασης είναι $\Delta f_{\mathrm{ss}} = -\Delta P / \lambda$, όπου $\lambda$ (MW/Hz) η
χαρακτηριστική ισχύος-συχνότητας της περιοχής, δηλαδή το άθροισμα της συνεισφοράς των
ρυθμιστών και της αυτορρύθμισης του φορτίου. Η ΕΔΣ επομένως σταθεροποιεί τη συχνότητα σε
νέα τιμή μόνιμης κατάστασης, με μόνιμο σφάλμα ως προς την ονομαστική τιμή, το οποίο
εξαλείφεται από την εφεδρεία αποκατάστασης συχνότητας.

\section{Εφεδρεία αποκατάστασης συχνότητας σε απομονωμένο σύστημα}
\label{sec:frr-single}

Η ΕΑΣ εξετάζεται αρχικά για \emph{απομονωμένο} σύστημα, χωρίς διασυνδέσεις με γειτονικά
συστήματα. Στην περίπτωση αυτή το μόνο μέγεθος ελέγχου είναι η απόκλιση συχνότητας και ο
κεντρικός ελεγκτής εφαρμόζει \emph{ολοκληρωτική} δράση:

\begin{equation}
\Delta P_{\mathrm{FRR}}(t) = -K_I \int_0^{t} \Delta f(\tau)\,\mathrm{d}\tau
\label{eq:frr-isolated}
\end{equation}

Όσο $\Delta f \neq 0$, το ολοκλήρωμα μεταβάλλεται και η εντολή παραγωγής αυξάνεται· σε μόνιμη
κατάσταση ισχύει κατ' ανάγκη $\Delta f = 0$. Ο ολοκληρωτικός έλεγχος εξαλείφει επομένως το
μόνιμο σφάλμα της αναλογικής ΕΔΣ. Με την επαναφορά της συχνότητας στην ονομαστική τιμή, η
συνεισφορά της ΕΔΣ μηδενίζεται σύμφωνα με την \eqref{eq:droop}, δηλαδή η ΕΑΣ
\emph{αποδεσμεύει} την ΕΔΣ, η οποία καθίσταται εκ νέου διαθέσιμη για επόμενη διαταραχή.

Η απόκριση της ΕΑΣ σχεδιάζεται σημαντικά βραδύτερη από εκείνη της ΕΔΣ (λεπτά έναντι
δευτερολέπτων). Ο διαχωρισμός των χρονικών κλιμάκων αποτρέπει την αλληλεπίδραση των δύο
βρόχων ελέγχου και τις ταλαντώσεις που θα προέκυπταν από αυτήν.

\section{Εφεδρεία αποκατάστασης συχνότητας σε σύστημα πολλών περιοχών}
\label{sec:frr-multi}

Σε διασυνδεδεμένο σύστημα η συχνότητα είναι κοινή σε ολόκληρη τη σύγχρονη περιοχή,
οπότε απώλεια παραγωγής σε οποιαδήποτε περιοχή προκαλεί απόκλιση συχνότητας σε όλες. Αν ο
ελεγκτής κάθε περιοχής χρησιμοποιούσε μόνο την απόκλιση συχνότητας, όλες οι περιοχές θα
ενεργοποιούσαν ΕΑΣ για την ίδια διαταραχή. Αυτό θα προκαλούσε υπερδιόρθωση, αποκλίσεις των
ανταλλαγών ισχύος από τα προγράμματα και μετακύλιση του κόστους εξισορρόπησης στις περιοχές
που δεν προκάλεσαν τη διαταραχή.

Για τον διαχωρισμό της περιοχής στην οποία εμφανίστηκε η ανισορροπία, ο ελεγκτής κάθε
περιοχής ελέγχου φορτίου-συχνότητας (περιοχή ΕΦΣ· Load-Frequency Control area, LFC area) μηδενίζει το \emph{σφάλμα ελέγχου
περιοχής} (ΣΕΠ· Area Control Error, ACE):

\begin{equation}
\mathrm{ACE} = \Delta P_{\mathrm{tie}} + K \, \Delta f
\label{eq:ace}
\end{equation}

όπου $\Delta P_{\mathrm{tie}}$ το σφάλμα ελέγχου ισχύος, δηλαδή η απόκλιση της μετρούμενης
ανταλλαγής ισχύος στις γραμμές διασύνδεσης από το πρόγραμμα, και $K$ ο συντελεστής K της
περιοχής σε MW/Hz (στην κλασική βιβλιογραφία συντελεστής πόλωσης συχνότητας, frequency
bias). Στο απομονωμένο σύστημα ο πρώτος όρος δεν υπάρχει και η \eqref{eq:ace} ανάγεται στο
κριτήριο της προηγούμενης ενότητας. Για τον λόγο αυτόν ο SOGL ορίζει ως ενιαίο μέγεθος
ελέγχου της αποκατάστασης το \emph{σφάλμα ελέγχου αποκατάστασης συχνότητας} (ΣΕΑΣ·
Frequency Restoration Control Error, FRCE): σε σύστημα πολλών περιοχών ισούται με το ΣΕΠ
της περιοχής, ενώ όταν η περιοχή ελέγχου συμπίπτει με ολόκληρη τη σύγχρονη περιοχή
ισούται με την απόκλιση συχνότητας.

Ο συνδυασμός των δύο όρων διακρίνει τη θέση της διαταραχής (Σχήμα~\ref{fig:lfc}):

\begin{itemize}
\item \textbf{Διαταραχή εντός της περιοχής.} Η συχνότητα μειώνεται ($\Delta f < 0$) και οι
      γειτονικές περιοχές παρέχουν ισχύ υποστήριξης, οπότε οι εισαγωγές υπερβαίνουν το
      πρόγραμμα ($\Delta P_{\mathrm{tie}} < 0$). Οι δύο όροι είναι ομόσημοι, το ΣΕΠ λαμβάνει
      μεγάλη απόλυτη τιμή και ενεργοποιείται ΕΑΣ στην περιοχή.
\item \textbf{Διαταραχή εκτός της περιοχής.} Η συχνότητα μειώνεται, αλλά η περιοχή εξάγει
      πρόσθετη ισχύ προς τη γειτονική περιοχή ($\Delta P_{\mathrm{tie}} > 0$). Για κατάλληλη
      τιμή του $K$ οι δύο όροι αλληλοαναιρούνται και το ΣΕΠ παραμένει σχεδόν μηδενικό· η
      περιοχή συμμετέχει μόνο μέσω της ΕΔΣ.
\end{itemize}

Η ιδιότητα αυτή ονομάζεται \emph{μη αλληλεπιδραστικός έλεγχος} (non-interactive control):
κάθε περιοχή αντισταθμίζει μόνο τις δικές της ανισορροπίες, ενώ όλες οι περιοχές
συμμετέχουν μέσω της ΕΔΣ στη συγκράτηση της συχνότητας αμέσως μετά τη διαταραχή. Η
αλληλοαναίρεση είναι πλήρης όταν ο συντελεστής $K$ ισούται με τη χαρακτηριστική
ισχύος-συχνότητας της περιοχής, δηλαδή με το άθροισμα της συνεισφοράς της ΕΔΣ, του
αυτοελέγχου της παραγωγής και της αυτορρύθμισης του φορτίου ανά Hz.

\begin{figure}[!htbp]
\centering
\includegraphics[width=\linewidth]{figures/lfc-areas.svg}
\caption{Εφεδρεία αποκατάστασης συχνότητας σε απομονωμένο και σε διασυνδεδεμένο σύστημα. Στην
πρώτη περίπτωση αρκεί η απόκλιση συχνότητας· στη δεύτερη απαιτείται και ο όρος των
ανταλλαγών ισχύος στις γραμμές διασύνδεσης, ώστε να ενεργοποιείται ΕΑΣ μόνο στην περιοχή
όπου εμφανίστηκε η διαταραχή.}
\label{fig:lfc}
\end{figure}

\section{Παράμετροι συχνότητας για την Ηπειρωτική Ευρώπη}
\label{sec:ce-params}

\begin{itemize}
\item Το \textbf{συμβάν αναφοράς} (reference incident) για τη διαστασιολόγηση της ΕΔΣ είναι
      ανισορροπία 3\,000 MW, είτε ως απώλεια παραγωγής είτε ως απώλεια φορτίου. Η ελάχιστη
      διαθέσιμη ΕΔΣ στη σύγχρονη περιοχή είναι ανά πάσα στιγμή 3\,000 MW.
\item Η \textbf{μέγιστη απόκλιση συχνότητας σταθερής κατάστασης} για το συμβάν αναφοράς
      είναι 200 mHz. Από τα δύο μεγέθη προκύπτει ο συντελεστής αναλογίας μεταξύ ΕΔΣ και
      απόκλισης συχνότητας: 3\,000 MW / 0,2 Hz = 15\,000 MW/Hz.
\item Η \textbf{τυπική περιοχή συχνότητας} (standard frequency range) είναι $\pm 50$ mHz· η
      συχνότητα δεν επιτρέπεται να βρίσκεται εκτός αυτής για περισσότερα από
      15\,000 λεπτά ανά έτος. Η μέγιστη στιγμιαία απόκλιση συχνότητας είναι 800 mHz.
\item Απόκλιση συχνότητας σταθερής κατάστασης μεγαλύτερη των 200 mHz κατατάσσει το σύστημα
      σε κατάσταση \emph{έκτακτης ανάγκης} (Κεφάλαιο 1, ενότητα
      «\nameref{sec:states}»).
\end{itemize}

\section{Περιοχές συχνότητας και τάσης λειτουργίας}
\label{sec:f-ranges}

Οι εφεδρείες συχνότητας περιορίζουν την απόκλιση της συχνότητας μόνο εφόσον οι μονάδες
ηλεκτροπαραγωγής παραμένουν συνδεδεμένες κατά τη διάρκεια της διαταραχής. Αν οι μονάδες
αποσυνδέονταν με τις πρώτες αποκλίσεις, η ανισορροπία θα αυξανόταν και η συχνότητα θα
απομακρυνόταν ακόμη περισσότερο από την ονομαστική τιμή. Για τον λόγο αυτόν ο RfG ορίζει περιοχές συχνότητας εντός των οποίων οι μονάδες πρέπει να παραμένουν
συνδεδεμένες και να λειτουργούν τουλάχιστον για ορισμένο χρόνο (Πίνακας~\ref{tab:f-ranges}).

\begin{table}[!htbp]
\centering
\begin{tabular}{lll}
\hline
\textbf{Περιοχή συχνότητας (Hz)} & \textbf{Απόκλιση από τα 50 Hz (Hz)}
& \textbf{Ελάχιστη διάρκεια λειτουργίας} \\
\hline
$<$ 47,5 & $<$ −2,5 & δεν απαιτείται \\
47,5--48,5 & −2,5 έως −1,5 & 30 min \\
48,5--49,0 & −1,5 έως −1,0 & 30 min \\
49,0--51,0 & −1,0 έως +1,0 & απεριόριστη \\
51,0--51,5 & +1,0 έως +1,5 & 30 min \\
$>$ 51,5 & $>$ +1,5 & δεν απαιτείται \\
\hline
\end{tabular}
\caption{Περιοχές συχνότητας και ελάχιστη διάρκεια λειτουργίας των μονάδων χωρίς αποσύνδεση
στην Ηπειρωτική Ευρώπη, κατά τον RfG και τον Κώδικα Διαχείρισης του ΕΣΜΗΕ.}
\label{tab:f-ranges}
\end{table}

Μεταξύ 49,0 και 51,0 Hz η λειτουργία είναι απεριόριστη, ενώ στις ακραίες περιοχές απαιτείται
για περιορισμένο χρόνο. Για την περιοχή 51,0--51,5 Hz ο RfG ορίζει 30 min. Για τις δύο
περιοχές υποσυχνότητας θέτει κατώτατο όριο 30 min και αφήνει την τελική τιμή στον ΔΣΜΗΕ· ο
Κώδικας Διαχείρισης του ΕΣΜΗΕ εφαρμόζει τα 30 min και στις δύο. Εκτός της περιοχής
47,5--51,5 Hz η παραμονή σε λειτουργία δεν απαιτείται και οι μονάδες μπορούν να αποσυνδεθούν
από τις προστασίες υποσυχνότητας και υπερσυχνότητας. Οι προστασίες αυτές προφυλάσσουν τον
εξοπλισμό, π.χ. τα πτερύγια των ατμοστροβίλων, τα οποία καταπονούνται σε παρατεταμένη
λειτουργία μακριά από την ονομαστική συχνότητα.

Οι περιοχές του Πίνακα~\ref{tab:f-ranges} είναι πολύ ευρύτερες από τις αποκλίσεις που
προκαλεί το συμβάν αναφοράς, με μέγιστη στιγμιαία απόκλιση 800 mHz (ενότητα
«\nameref{sec:ce-params}»). Αφορούν σοβαρότερες διαταραχές, όπως ο διαχωρισμός της
σύγχρονης περιοχής σε τμήματα με μεγάλο έλλειμμα ή πλεόνασμα παραγωγής. Τότε η πτώση της
συχνότητας ανακόπτεται από την αυτόματη αποσύνδεση ζήτησης λόγω χαμηλής συχνότητας του
σχεδίου άμυνας του συστήματος, η οποία στην Ηπειρωτική Ευρώπη ξεκινά στα 49,0 Hz και
ολοκληρώνεται κλιμακωτά έως τα 48,0 Hz. Η αποσύνδεση ζήτησης αποκαθιστά την ισορροπία μόνο
αν οι μονάδες παραμένουν συνδεδεμένες σε όλη αυτή την περιοχή· γι' αυτό το κατώτατο όριο των
47,5 Hz βρίσκεται κάτω από το τελευταίο της βήμα.

Οι απαιτήσεις για τη συχνότητα συνδυάζονται με αντίστοιχες απαιτήσεις για την τάση στο
σημείο σύνδεσης. Για τις μονάδες που συνδέονται στο σύστημα μεταφοράς (τύπου Δ) ο RfG ορίζει
απεριόριστη λειτουργία στις περιοχές 0,90--1,118 p.u. στα 150 kV και 0,90--1,05 p.u. στα
400 kV, οι οποίες συμπίπτουν με τα όρια κανονικής λειτουργίας των ζυγών του ΕΣΜΗΕ
(Κεφάλαιο 2, ενότητα «\nameref{sec:v-control}»). Στη ζώνη 0,85--0,90 p.u. απαιτεί λειτουργία
για 60 min, ενώ στη ζώνη υπέρτασης, έως 1,15 p.u. στα 150 kV και έως 1,10 p.u. στα 400 kV,
για χρόνο μεταξύ 20 και 60 min, τον οποίο ορίζει ο ΔΣΜΗΕ.

Ο Κώδικας Διαχείρισης του ΕΣΜΗΕ συνδυάζει τις δύο απαιτήσεις σε διαγράμματα
συχνότητας-τάσης, π.χ. στις απαιτήσεις για τους σταθμούς αποθήκευσης ηλεκτρικής ενέργειας που
συνδέονται στο σύστημα μεταφοράς, όπου και η ζώνη υπέρτασης λαμβάνει 60 min. Το
Σχήμα~\ref{fig:f-u} αποδίδει τα διαγράμματα για σημείο σύνδεσης στα 150 kV και στα 400 kV.
Σε κάθε σημείο λειτουργίας $(f, U)$ ισχύει ο μικρότερος από τους χρόνους που προκύπτουν
χωριστά για τη συχνότητα και για την τάση. Η λειτουργία είναι επομένως απεριόριστη μόνο όταν
και τα δύο μεγέθη βρίσκονται εντός των ορίων κανονικής λειτουργίας. Περιορίζεται σε 60 min
στις ζώνες μικρής απόκλισης της τάσης και σε 30 min, σε όλο το εύρος τάσης, εκτός της
περιοχής 49,0--51,0 Hz.

\begin{figure}[!htbp]
\centering
\includegraphics[width=\linewidth]{figures/freq-voltage-ranges.svg}
\caption{Διαγράμματα συχνότητας-τάσης για σημείο σύνδεσης στα 150 kV (αριστερά) και στα
400 kV (δεξιά), κατά τον Κώδικα Διαχείρισης του ΕΣΜΗΕ: ελάχιστη διάρκεια λειτουργίας σε
σταθερή κατάσταση για κάθε συνδυασμό της συχνότητας $f$ του συστήματος και της τάσης $U$ στο
σημείο σύνδεσης. Εκτός των σκιασμένων περιοχών η παραμονή σε λειτουργία δεν απαιτείται.}
\label{fig:f-u}
\end{figure}

Στις γωνίες των διαγραμμάτων οι δύο αποκλίσεις συνδυάζονται δυσμενώς για τις σύγχρονες
γεννήτριες και τους μετασχηματιστές. Η μαγνητική ροή στον πυρήνα τους είναι ανάλογη του λόγου
$U/f$, ο οποίος αυξάνεται με ταυτόχρονη υπέρταση και υποσυχνότητα: στα 47,5~Hz (0,95~p.u.) και
με τάση 1,10~p.u. ο λόγος φθάνει τα 1,16 p.u. Όταν υπερβαίνει την ονομαστική τιμή κατά
περισσότερο από περίπου 5\%, ο πυρήνας οδηγείται σε κορεσμό, με αυξημένες απώλειες και
υπερθέρμανση (υπερδιέγερση). Για τον λόγο αυτόν ο RfG επιτρέπει στον ΔΣΜΗΕ να ορίζει
μικρότερους χρόνους λειτουργίας για ταυτόχρονη υπέρταση και υποσυχνότητα ή ταυτόχρονη
υπόταση και υπερσυχνότητα, ενώ οι σύγχρονες γεννήτριες διαθέτουν προστασία υπερδιέγερσης
(V/Hz).

Τα διαγράμματα αφορούν λειτουργία σε σταθερή κατάσταση. Οι βυθίσεις της τάσης κάτω από
0,85 p.u. κατά τη διάρκεια σφαλμάτων καλύπτονται από τις απαιτήσεις αδιάλειπτης λειτουργίας
υπό σφάλμα (Κεφάλαιο 2, υποενότητα «\nameref{sec:frt-req}»), ενώ οι γρήγορες μεταβολές της
συχνότητας από την απαίτηση αντοχής στον RoCoF (ενότητα
«\nameref{sec:inertia}»), την οποία ορίζει επίσης ο RfG.

% TODO: παραδείγματα και λυμένες ασκήσεις (σελίδα 03b-paradeigmata.md, όπως στο Κεφάλαιο 1)
